\section{Conclusion}
\label{sec:conclusion}
% \cnnote{结论待写}{正式结果尚未汇齐。完成主对比、五组消融、数据质量及评价核验后，再根据实际证据写结论；当前不新增性能、泛化或交互体验提升的结论。}
\paragraph{Limitations and future work.}
% \cnnote{局限段落待写}{结果确定后，用一个短段落结合失败例说明适用边界。当前设定只使用当前语音与首帧；自动标注、外部事件检测及单一参考评价的局限需结合核验结果说明。这里不新增多轮历史、人物长期一致性或实时流式实验要求。}
% \cnnote{作者自查}{贡献：保留数据、模型和评价三条方向，效果证据待补。写作：已区分标注／生成模型、训练／推理描述和两类事件协议。实验：五组已规划，成绩留空。评价：外部检测、人工规范和覆盖待落实。方法：冻结基座及梯度路径已说明，版本与细节待核对。完整逐项核对见review/含义与证据核对.md。}
